\documentclass[11pt]{amsart}
\usepackage{geometry}                % See geometry.pdf to learn the layout options. There are lots.
\geometry{letterpaper}                   % ... or a4paper or a5paper or ... 
%\geometry{landscape}                % Activate for for rotated page geometry
%\usepackage[parfill]{parskip}    % Activate to begin paragraphs with an empty line rather than an indent
\usepackage{graphicx}
\usepackage{amssymb}
\usepackage{epstopdf}
\DeclareGraphicsRule{.tif}{png}{.png}{`convert #1 `dirname #1`/`basename #1 .tif`.png}

\title{PS 3.7}
%\date{}                                           % Activate to display a given date or no date

\begin{document}
\maketitle
%\section{}
%\subsection{}
\noindent Calcule el aumento de sueldos para N empleados de una empresa, bajo el si­guiente criterio:
\newline

• Si el sueldo es menor a \$10,000 :  Aumento 10\%

• Si el sueldo está comprendido entre \$10,000 y \$25,000 :  Aumento 7\% 

• Si el sueldo es mayor a \$25,000 : Aumento 8\%
\newline

\noindent Imprima lo siguiente:
\newline

a) El nuevo sueldo del trabajador

b) El monto total de la nómina considerando el aumento.
\newline

\noindent Datos: N, $SUE_{1}$, $SUE_{2}$, . . $SUE_{N}$
\newline

\noindent Donde:
\newline

N es una variable de tipo entero que representa el número de em­pleados de la empresa.
\newline

SUE¡. es una variable de tipo real que representa el sueldo del trabajador i (1 $\leq$ i $\leq$ N).


\end{document}  